    \title{Lista de Exercícios: IDS/IPS}
        \author{Prof. Gabriel Rodrigues Caldas de Aquino}
        \date{Compilado em: \\ \today}

\begin{document}

\maketitle

\section{Detecção e Prevenção de Intrusões (IDS/IPS)}
Os sistemas de detecção e prevenção de intrusões representam uma camada  na estratégia de defesa em profundidade, complementando as defesas de perímetro.

\begin{enumerate}
    \item Defina o que é uma intrusão
    \item Defina o que é um \textbf{IDS (Sistema de Detecção de Intrusão)}.
    \item Defina o que é um \textbf{IPS (Sistema de Prevenção de Intrusão)}.
    \item Diga qual é a diferença fundamental entre eles em termos de ação tomada quando uma intrusão é identificada?
    \item Diga como que o uso de um firewall do tipo \textbf{filtro de pacotes} pode ser insuficiente para o trabalho de detecção de intrusão.
\end{enumerate}

\section{Detectando uma Intrusão}
Para detectar uma intrusão é necessário analisar se o comportamento observado representa uma ameaça. Considerando que existem dois tipos fundamentais de comportamentos: (i) aqueles que já conhecemos e temos catalogados e (ii) aqueles que ainda não conhecemos.

\begin{enumerate}
    \item \textbf{Defina} e \textbf{contraste} as duas principais abordagens de análise em IDS para lidar com esses diferentes tipos de comportamentos:
          \begin{itemize}
              \item \textbf{Detecção por Assinaturas}
              \item \textbf{Detecção de Anomalia}
          \end{itemize}
    \item \textbf{Explique} qual abordagem é mais eficaz para cada tipo de comportamento mencionado e \textbf{por quê}.
    \item \textbf{Analise} as diferenças entre Detecção por Assinatura e Detecção de Anomalia em termos de:
          \begin{itemize}
              \item Capacidade de detectar ataques zero-day
              \item Taxa de falsos positivos
              \item Custo computacional e de manutenção
              \item Adaptabilidade a novas ameaças
          \end{itemize}

\end{enumerate}

\section{NIDS e HIDS}
Considere as arquiteturas de sistemas de detecção de intrusão baseadas em host e em rede:
\begin{enumerate}
    \item  Defina e contraponha as características de:
          \begin{itemize}
              \item     HIDS (Host-based Intrusion Detection System)

              \item   NIDS (Network-based Intrusion Detection System)
          \end{itemize}
    \item  Explique para cada um dos seguintes cenários qual abordagem (HIDS ou NIDS) seria mais adequada e por quê
          \begin{itemize}
              \item     Detecção de um malware que realiza chamadas de sistema anômalas

              \item       Identificação de varredura de portas na rede corporativa

              \item        Detecção de tentativas de exploração de vulnerabilidade em um servidor web

              \item       Monitoramento de acesso não autorizado a arquivos sensíveis
          \end{itemize}
    \item Nos cenários citados acima, é possível que um uso em conjunto de NIDS e HIDS pode melhorar a detecção? Justifique como o uso poderia melhorar ou não a detecção em cada cenário.
    \item  Descreva o passo a passo de funcionamento do SNORT como um NIDS, desde a captura do pacote até a geração do alerta.
    \item Analise uma limitação significativa do NIDS em ambientes com tráfego criptografado e proponha uma estratégia complementar para mitigar esta limitação.
\end{enumerate}


\section{Regra de detecção - Snort}
Considere a seguinte regra do SNORT para detecção de tráfego BitTorrent:

\begin{verbatim}
alert tcp $HOME_NET any -> $EXTERNAL_NET ![7680,1521] 
(msg:"ET P2P BitTorrent peer sync"; 
 flow:established,to_server; 
 content:"|00 00 00 0d 06 00|"; depth:6; 
 threshold: type limit, track by_dst, seconds 300, count 1; 
 reference:url,bitconjurer.org/BitTorrent/protocol.html; 
 classtype:policy-violation; sid:2000334; rev:14; 
 metadata:created_at 2010_07_30, confidence Medium, 
 signature_severity Informational, updated_at 2025_06_30;)
\end{verbatim}

\begin{enumerate}
    \item Descreva o que cada parte do cabeçalho da regra especifica:
          \begin{itemize}
              \item     Ação a ser tomada
              \item    Protocolo e direção do tráfego
              \item     Endereços de origem e destino
              \item     Portas envolvidas
          \end{itemize}
    \item  Explique a função de cada uma das seguintes opções:
          \begin{itemize}
              \item   \texttt{flow:established,to\_server}
              \item   \texttt{content:"|00 00 00 0d 06 00|"; depth:6}
              \item   \texttt{threshold: type limit, track by\_dst, seconds 300, count 1}
          \end{itemize}
    \item Identifique que tipo de abordagem de detecção esta regra representa (se é por Assinatura ou Anomalia) e justifique sua resposta.
\end{enumerate}

\end{document}